\documentclass[11pt]{article}
\usepackage[margin=1in]{geometry}
\usepackage{hyperref}
\usepackage{enumitem}
\title{Intro Typing and Simple Circuits}
\author{Piter Garcia}
\date{\today}
\begin{document}
\maketitle

\section*{Grade and Class Match}
\textbf{Grade:} Kindergarten\\
\textbf{Likely blocks:} Day 1 · 12:25-1:15 · Justinger; Day 2 · 12:25-1:15 · Kesys; Day 4 · 12:25-1:15 · Pum

\section*{Lesson Focus}
Short keyboarding practice paired with an early introduction to circuits and safe materials handling.

\section*{Learning Targets}
\begin{itemize}[leftmargin=*]
\item Students practice simple keyboard or touch-entry routines.
\item Students identify a battery, light, or connection in a circuit demonstration.
\item Students follow material expectations safely in a guided activity.
\end{itemize}

\section*{Materials}
\begin{itemize}[leftmargin=*]
\item Student devices
\item Typing practice activity
\item Circuit demo materials
\end{itemize}

\section*{Suggested Lesson Flow}
\begin{enumerate}[leftmargin=*]
\item Open with the exact device, tool, or routine students need in order to start successfully.
\item Model one example task before students begin independent or partner work.
\item Give students time to build, code, design, or document their work while circulating for support.
\item Close with a short share-out, reflection, or debugging discussion.
\end{enumerate}

\section*{Source Notes}
This lesson packet is enriched with transcript-derived lesson notes, the school schedule roster, and official starting-point resources. See the paired Markdown guide for clickable links.

\bibliographystyle{plain}
\bibliography{intro-typing-and-simple-circuits}
\end{document}
